% =============================================================================
% P4: Evaluación de la Inversión de Dependencias en la Práctica: OrderProcessor
% =============================================================================
% Fragmento puro LaTeX — se incluye con \include desde SemArq1_TeamBales.tex
% Prohibido: \documentclass, \begin{document}, \usepackage, preámbulo.
% =============================================================================

\chapter{P4: Evaluación de la Inversión de Dependencias en la Práctica: OrderProcessor}

\section{a) ¿Por qué viola la dependencia directa los principios de Hexagonal, Limpia y Onion?}

Un equipo de desarrollo trabaja en un sistema de gestión de pedidos. La lógica de negocio principal reside en la clase \texttt{OrderProcessor}, que necesita comunicarse con un servicio externo de validación de direcciones. La implementación actual tiene una dependencia directa:

\begin{center}
\begin{tikzpicture}[font=\sffamily\small,
    clase/.style={modulo, minimum width=3.9cm, font=\ttfamily\footnotesize, fill=blue!5, draw=blue!50!black},
    ext/.style={modulo, minimum width=3.4cm, font=\ttfamily\footnotesize, fill=gray!10, draw=gray!60!black}]
  \node[clase] (op)  at (0,0)    {OrderProcessor};
  \node[clase] (avs) at (6.2,0)  {AddressValidatorService};
  \node[ext]   (lib) at (11.9,0) {HttpLib};
  \draw[dependencia] (op)  -- node[above, font=\sffamily\scriptsize] {instancia y llama} (avs);
  \draw[dependencia] (avs) -- node[above, font=\sffamily\scriptsize] {usa} (lib);
\end{tikzpicture}
\end{center}

Esta arquitectura viola los principios fundamentales de las arquitecturas Hexagonal, Limpia y Onion por las siguientes razones:

\subsection*{Violación del Principio de Inversión de Dependencias (DIP)}

El DIP establece que \enquote{las dependencias del código fuente deben referirse solo a abstracciones, no a concreciones} \fuente{martin2017}{87}. En la implementación actual:

\begin{enumerate}
    \item \texttt{OrderProcessor}, la política de alto nivel, depende directamente de un detalle de implementación concreto y volátil: \texttt{AddressValidatorService}.
    \item El servicio externo es inherentemente volátil: puede cambiar su API, su contrato, su disponibilidad, o ser reemplazado por otro proveedor.
    \item Como señala \citet[p.~88]{martin2017}: \enquote{Every change to an abstract interface corresponds to a change to its concrete implementations. Conversely, changes to concrete implementations do not always, or even usually, require changes to the interfaces that they implement. Therefore interfaces are less volatile than implementations.}
\end{enumerate}

\subsection*{Violación por arquitectura}

\begin{itemize}
    \item \textbf{Arquitectura Hexagonal:} La lógica de negocio (interior del hexágono) depende de un servicio externo (adaptador), invirtiendo la relación correcta. Los puertos (interfaces) deben definirse en el interior y los adaptadores (implementaciones) en el exterior; los adaptadores dependen de los puertos, no al revés \citep[«Nature of the Solution»]{cockburn2005}.
    \item \textbf{Arquitectura Limpia:} Se viola la Regla de Dependencia: \enquote{Source code dependencies must point only inward, toward higher-level policies. Nothing in an inner circle can know anything at all about something in an outer circle} \fuente{martin2017}{203}. Las dependencias apuntan hacia afuera (hacia detalles volátiles) en lugar de hacia adentro (hacia políticas de alto nivel).
    \item \textbf{Arquitectura de Cebolla:} La dependencia se aleja del centro (Modelo de Dominio). En Onion todas las dependencias deben apuntar hacia el centro, y las capas externas deben implementar interfaces definidas en las capas internas \fuente{martin2017}{202}.
\end{itemize}

\subsection*{Acoplamiento transitivo a la librería HTTP}

La cadena \texttt{OrderProcessor} $\rightarrow$ \texttt{AddressValidatorService} $\rightarrow$ \texttt{HttpClient} acopla transitivamente la lógica de negocio con una librería HTTP concreta. Cualquier cambio en la API del servicio, su contrato o la librería HTTP obliga a modificar \texttt{OrderProcessor}, la clase que aloja la política de negocio principal.

\subsection*{Dirección correcta de las dependencias}

El DIP y la Regla de Dependencia exigen que las dependencias del código fuente apunten hacia las abstracciones del núcleo:

\begin{center}
\begin{tikzpicture}[font=\sffamily\small,
    clase/.style={modulo, minimum width=3.9cm, font=\ttfamily\footnotesize, fill=blue!5, draw=blue!50!black},
    puerto/.style={interfaz, minimum width=3.6cm, font=\ttfamily\itshape\footnotesize, draw=green!40!black, fill=green!8},
    implementa/.style={-{Triangle[open, length=3mm, width=3mm]}, thick, dashed, shorten >=1pt, shorten <=1pt}]
  \node[clase]  (op)  at (0,0)    {OrderProcessor};
  \node[puerto] (iav) at (5.9,0)  {IAddressValidator};
  \node[clase]  (avs) at (11.6,0) {HttpAddressValidatorAdapter};
  \draw[dependencia] (op)  -- node[above, font=\sffamily\scriptsize] {usa} (iav);
  \draw[implementa]  (avs) -- node[above, font=\sffamily\scriptsize] {implementa} (iav);
\end{tikzpicture}
\end{center}

La política de alto nivel y el detalle dependen ahora de la misma abstracción, definida en el núcleo: \texttt{OrderProcessor} usa \texttt{IAddressValidator} y \texttt{AddressValidatorService} la implementa. Nótese que el \textit{flujo de control} se mantiene —\texttt{OrderProcessor} sigue solicitando la validación—; lo que se invierte es la \textit{dependencia del código fuente}, que ahora apunta desde el detalle hacia la abstracción \fuente{martin2017}{87, 91}.

\section{b) Rediseño con inversión de dependencias}

\subsection*{Propuesta de solución}

Para alinear el diseño con los principios de Clean Architecture, Arquitectura Hexagonal y Arquitectura de Cebolla —arquitecturas cuya característica común es producir sistemas independientes de frameworks, UI, base de datos y agentes externos \fuente{martin2017}{202}, regidas por la Regla de Dependencia \fuente{martin2017}{203}— se realizan tres cambios:

\begin{enumerate}
    \item \textbf{Crear una abstracción en el núcleo (Capa de Dominio/Entidades):}
    \begin{itemize}
        \item Definir la interfaz \texttt{IAddressValidator} (puerto de validación) en la capa más interna.
        \item Define el contrato: \texttt{validarDireccion(Direccion): ResultadoValidacion}.
        \item La interfaz es estable y poco volátil: \enquote{Don't refer to volatile concrete classes. Refer to abstract interfaces instead} \fuente{martin2017}{89}.
    \end{itemize}

    \item \textbf{Implementar la interfaz en la capa de infraestructura (exterior del hexágono):}
    \begin{itemize}
        \item \texttt{AddressValidatorService} implementa \texttt{IAddressValidator}.
        \item Reside en la capa de infraestructura (adaptador externo en Hexagonal).
        \item Encapsula la dependencia de la librería HTTP concreta.
    \end{itemize}

    \item \textbf{Inyectar la dependencia en \texttt{OrderProcessor}:}
    \begin{itemize}
        \item \texttt{OrderProcessor} recibe \texttt{IAddressValidator} por constructor (inyección de dependencias).
        \item \texttt{OrderProcessor} \textbf{nunca} instancia directamente \texttt{AddressValidatorService}.
        \item Se usa \textit{Abstract Factory} o contenedor de DI para la creación \fuente{martin2017}{89}.
    \end{itemize}
\end{enumerate}

\subsection*{Diagrama de dependencias antes/después}

\begin{figure}[H]
    \centering
    \begin{tikzpicture}[font=\sffamily\small,
        clase/.style={modulo, minimum width=3.9cm, minimum height=1.05cm, font=\ttfamily\footnotesize, fill=blue!5, draw=blue!50!black},
        puerto/.style={interfaz, minimum width=3.9cm, minimum height=1.05cm, font=\ttfamily\itshape\footnotesize, draw=green!40!black, fill=green!8},
        ext/.style={modulo, minimum width=3.4cm, minimum height=1.05cm, font=\ttfamily\footnotesize, fill=gray!10, draw=gray!60!black},
        implementa/.style={-{Triangle[open, length=3mm, width=3mm]}, thick, dashed, shorten >=1pt, shorten <=1pt},
        etiqueta/.style={font=\sffamily\scriptsize, fill=white, inner sep=1.5pt},
        titulo/.style={font=\sffamily\small\bfseries, anchor=west},
        zona/.style={rounded corners=6pt, thick},
        nombrezona/.style={font=\sffamily\scriptsize\bfseries, anchor=north west}]

        % ---------- (a) ANTES ----------
        \node[titulo, red!60!black] at (-2.1,1.35) {(a) Antes: la política depende del detalle};
        \node[clase] (op1)  at (0,0)    {OrderProcessor};
        \node[clase] (avs1) at (6.2,0)  {AddressValidatorService};
        \node[ext]   (lib1) at (11.9,0) {HttpLib};
        \draw[dependencia, red!60!black] (op1)  -- node[etiqueta, above=2pt] {instancia y llama} (avs1);
        \draw[dependencia] (avs1) -- node[etiqueta, above=2pt] {usa} (lib1);

        % ---------- (b) DESPUÉS ----------
        \node[titulo, green!35!black] at (-2.1,-1.65) {(b) Después: el detalle depende de la abstracción};
        % Zonas
        \draw[zona, draw=blue!40!black, fill=blue!3]  (-2.25,-2.15) rectangle (2.25,-6.35);
        \node[nombrezona, blue!40!black] at (-2.15,-2.2) {NÚCLEO (dominio)};
        \draw[zona, draw=gray!55!black, fill=gray!4]   (3.6,-4.2) rectangle (14.0,-6.35);
        \node[nombrezona, gray!40!black] at (3.6,-4.25) {INFRAESTRUCTURA (exterior del hexágono)};
        % Nodos
        \node[clase]  (op2)  at (0,-3.35)   {OrderProcessor};
        \node[puerto] (iav)  at (0,-5.45)   {IAddressValidator\\[-1pt]{\sffamily\upshape\scriptsize puerto}};
        \node[clase]  (avs2) at (6.4,-5.45) {HttpAddressValidatorAdapter\\[-1pt]{\sffamily\scriptsize adaptador}};
        \node[ext]    (lib2) at (12.1,-5.45) {HttpLib};
        % Dependencias
        \draw[dependencia] (op2) -- node[etiqueta, right=2pt] {usa} (iav);
        \draw[implementa]  (avs2) -- node[etiqueta, above=2pt] {implementa} (iav);
        \draw[dependencia] (avs2) -- node[etiqueta, above=2pt] {usa} (lib2);
    \end{tikzpicture}
    \caption{Dependencias del código fuente antes y después de la inversión. En (a) la lógica de negocio depende de un detalle concreto; en (b) tanto \texttt{OrderProcessor} como el adaptador dependen de la interfaz definida en el núcleo, y la única dependencia que cruza la frontera apunta hacia dentro. \texttt{HttpAddressValidatorAdapter} es el \texttt{AddressValidatorService} original rediseñado como adaptador.}
    \label{fig:dip-before-after}
\end{figure}

\subsection*{Justificación teórica}

\begin{itemize}
    \item \textbf{Regla de Dependencia} \fuente{martin2017}{203}: Las dependencias apuntan hacia adentro; \texttt{AddressValidatorService} (infraestructura) depende de \texttt{IAddressValidator} (núcleo).
    \item \textbf{Independencia de frameworks} \fuente{martin2017}{202}: La arquitectura no depende de librerías HTTP; se usan como herramientas intercambiables.
    \item \textbf{Testabilidad} \fuente{martin2017}{202}: El núcleo se prueba sin UI, BD, servidor web ni servicios externos.
    \item \textbf{Propiedad del contrato}: \texttt{OrderProcessor} define \textit{qué} necesita (el contrato \texttt{IAddressValidator}); la infraestructura provee \textit{cómo} se cumple. El detalle depende de la abstracción del núcleo.
\end{itemize}

\section{c) Pruebas aisladas sin HTTP y mantenibilidad al cambiar de proveedor}

\subsection*{Aislamiento de la lógica de negocio}

Tras el rediseño, \texttt{OrderProcessor} no se encarga de la validación de direcciones ni de detalles de comunicación HTTP. Su única responsabilidad es la lógica de negocio de procesamiento de pedidos. La validación se delega a la implementación inyectada de \texttt{IAddressValidator}.

\subsection*{Pruebas unitarias puras con test doubles}

Ahora \texttt{OrderProcessor} es testeable de forma aislada, rápida y determinista:

\begin{enumerate}
    \item \textbf{Mock/Stub de \texttt{IAddressValidator}:} En las pruebas unitarias se inyecta un \textit{test double} que implementa \texttt{IAddressValidator} y devuelve respuestas controladas (válida, inválida, error de red, timeout, etc.).
    \item \textbf{Sin dependencias externas:} No se requiere conexión a red, servicio externo real, ni base de datos.
    \item \textbf{Cobertura de casos frontera:} Fácil simulación de errores de validación, timeouts, respuestas malformadas.
    \item \textbf{Velocidad:} Las pruebas se ejecutan en milisegundos, permitiendo integración continua rápida.
\end{enumerate}

\subsection*{Implicaciones para la mantenibilidad y el cambio de proveedor}

El diseño refactorizado confina la volatilidad a un único punto del sistema:

\begin{itemize}
    \item \textbf{Cambio de proveedor:} Sustituir el proveedor de validación solo requiere escribir un nuevo adaptador que implemente \texttt{IAddressValidator} y reconfigurar el mecanismo de construcción (Abstract Factory o contenedor DI); \texttt{OrderProcessor} y el resto del núcleo permanecen intactos.
    \item \textbf{Volatilidad contenida:} Dado que las interfaces son menos volátiles que las implementaciones \fuente{martin2017}{88}, el núcleo depende del elemento más estable; la volatilidad del proveedor queda confinada a la capa de infraestructura.
    \item \textbf{Extensión sin modificación:} Agregar un segundo proveedor (redundancia, tolerancia a fallos) no exige modificar el núcleo, solo sumar un nuevo adaptador que implemente la misma interfaz.
    \item \textbf{Impacto localizado de cambios:} Un cambio de contrato (interfaz) obliga a actualizar núcleo y adaptadores y se detecta en tiempo de compilación; un cambio de implementación no propaga modificaciones hacia el núcleo.
\end{itemize}

\subsection*{Principio subyacente}

Como establece \citet[p.~202]{martin2017}, una característica fundamental de estas arquitecturas es ser \textbf{Testable}: \enquote{The business rules can be tested without the UI, database, web server, or any other external element.} La inversión de dependencia es el mecanismo que hace esto posible.



\section*{Notas de investigación}

% =============================================================================
% Integrante: Enrique A Gonzalez Moreira
% Fragmento puro: se incluye desde preguntas/p4-inversion-dependencias/solucion.tex
% mediante % =============================================================================
% Integrante: Enrique A Gonzalez Moreira
% Fragmento puro: se incluye desde preguntas/p4-inversion-dependencias/solucion.tex
% mediante % =============================================================================
% Integrante: Enrique A Gonzalez Moreira
% Fragmento puro: se incluye desde preguntas/p4-inversion-dependencias/solucion.tex
% mediante \input{investigacion/p4-dip-interfaces/notas} (ya presente en ese
% archivo: sin este fragmento, SemArq1_TeamBales.tex no compila).
% Fuente de trabajo: notas.md (misma carpeta).
% Citas: páginas impresas de Martin (2017); Cockburn (2005) y Palermo (2008) no
% tienen paginación y se citan por sección o parte. Todas las claves están en
% core.bib (martin2017, cockburn2005, palermo2008): fuentes.bib no declara nada.
% Comillas: se usan comillas tipográficas literales, no \enquote, para no
% depender de csquotes (no está en la lista de paquetes de preamble.tex, #13).
% =============================================================================

% Grupo local: \emergencystretch evita líneas desbordadas sin afectar a otros capítulos.
\begingroup\emergencystretch=3em \hbadness=2000
\section*{Notas de investigación: DIP, Inversión de Dependencias y diseño de interfaces en Clean Architecture}

\section{Síntesis}

El Principio de Inversión de Dependencias (DIP) sostiene que los sistemas más flexibles son
aquellos en los que las dependencias del código fuente se refieren solo a abstracciones, no a
concreciones \parencite[87]{martin2017}. La consecuencia práctica es que una política de alto
nivel —una regla de negocio— no debe nombrar un detalle de bajo nivel: ni una base de datos, ni
un servicio externo, ni una librería de transporte. Ambos, política y detalle, dependen de una
abstracción declarada en el núcleo. Martin añade el criterio que convierte el principio en
decisión de diseño: toda modificación de una interfaz abstracta obliga a modificar sus
implementaciones, pero no al contrario, de modo que las interfaces son menos volátiles que las
implementaciones \parencite[88]{martin2017}; de ahí las reglas operativas de no referirse a
clases concretas volátiles, no sobreescribir funciones concretas y recurrir a factorías
abstractas para la creación de objetos \parencite[89]{martin2017}.

La Regla de Dependencia formaliza el principio a escala arquitectónica: las dependencias del
código fuente deben apuntar solo hacia adentro, hacia políticas de mayor nivel, y nada de un
círculo interno puede conocer nada de un círculo externo \parencite[203]{martin2017}. Martin
presenta esa regla como denominador común de un conjunto de arquitecturas —Hexagonal de
Cockburn, DCI, BCE— que persiguen la misma separación de intereses y producen sistemas
independientes de frameworks, de la interfaz de usuario, de la base de datos y de cualquier
agencia externa, y con reglas de negocio probables sin ninguno de esos elementos
\parencite[202]{martin2017}. En la Hexagonal, el núcleo declara los \emph{puertos} y la
tecnología aporta los \emph{adaptadores} que los implementan \parencite{cockburn2005}; en la
Onion, el Modelo de Dominio ocupa el centro y las interfaces de repositorio viven en el núcleo
mientras el comportamiento que toca la base de datos queda fuera \parencite{palermo2008}.

El mecanismo concreto es siempre el mismo: el detalle implementa la abstracción que el núcleo
declara, y la traducción entre el contrato del núcleo y la tecnología externa ocurre en la capa
de adaptadores de interfaz \parencite[205]{martin2017}. Conviene no confundirlo con la inversión
del flujo de control: el núcleo sigue llamando a la validación o a la consulta; lo que se
invierte es la dirección de la dependencia del código fuente, que en el límite cruzado pasa a
apuntar desde el detalle hacia la abstracción \parencite[87, 91]{martin2017}. Ese es el punto de
apoyo de la testabilidad: con el puerto en el núcleo, la prueba sustituye el adaptador real por
un doble controlado y ejercita la regla de negocio sin red, sin servidor y sin base de datos
\parencite[202]{martin2017}.

\section{Conceptos clave}

\begin{table}[htbp]
  \centering
  \footnotesize
  \renewcommand{\arraystretch}{1.2}
  \begin{tabular}{@{}p{0.22\linewidth}p{0.46\linewidth}p{0.24\linewidth}@{}}
    \toprule
    \textbf{Concepto} & \textbf{Definición} & \textbf{Fuente} \\
    \midrule
    \raggedright Principio de Inversión de Dependencias (DIP)
      & \raggedright Los sistemas más flexibles son los que dependen solo de abstracciones, no
        de concreciones; la política de alto nivel no nombra el detalle
      & \raggedright \textcite[87]{martin2017} \tabularnewline
    \raggedright Volatilidad de interfaces frente a implementaciones
      & \raggedright Cambiar una interfaz obliga a cambiar sus implementaciones, pero no al
        contrario: la interfaz es el elemento más estable del que depender
      & \raggedright \textcite[88]{martin2017} \tabularnewline
    \raggedright Reglas operativas del DIP
      & \raggedright No referirse a clases concretas volátiles, no sobreescribir funciones
        concretas y usar factorías abstractas para la creación
      & \raggedright \textcite[89]{martin2017} \tabularnewline
    \raggedright Regla de Dependencia
      & \raggedright Las dependencias del código fuente apuntan solo hacia adentro; nada de un
        círculo interno conoce nada de un círculo externo
      & \raggedright \textcite[203]{martin2017} \tabularnewline
    \raggedright Características comunes de Clean, Hexagonal, DCI y BCE
      & \raggedright Independencia de frameworks, de la interfaz de usuario, de la base de datos
        y de agencias externas; reglas de negocio probables en aislamiento
      & \raggedright \textcite[202]{martin2017} \tabularnewline
    \raggedright Puertos y adaptadores
      & \raggedright El núcleo declara el puerto y la tecnología aporta el adaptador que lo
        implementa, de modo que el núcleo ignora el dispositivo concreto
      & \raggedright \textcite{cockburn2005} \tabularnewline
    \raggedright Capa de adaptadores de interfaz
      & \raggedright Convierte los datos entre el formato del caso de uso y el de la tecnología
        externa; es el lugar natural del detalle
      & \raggedright \textcite[205]{martin2017} \tabularnewline
    \raggedright Dominio al centro (Onion)
      & \raggedright El Modelo de Dominio no se acopla más que a sí mismo; las interfaces de
        repositorio van en el núcleo y el acceso a datos fuera
      & \raggedright \textcite{palermo2008} \tabularnewline
    \raggedright Flujo de control frente a dependencia de código fuente
      & \raggedright El núcleo sigue invocando la operación; lo que se invierte en el límite es
        la dirección de la dependencia del código fuente
      & \raggedright \textcite[87, 91]{martin2017} \tabularnewline
    \raggedright Testabilidad mediante dobles de prueba
      & \raggedright Con el puerto en el núcleo, la prueba sustituye el adaptador real y ejercita
        la regla de negocio sin interfaz, base de datos ni servidor
      & \raggedright \textcite[202]{martin2017} \tabularnewline
    \bottomrule
  \end{tabular}
  \caption{Conceptos clave de la investigación INV-P4.}
  \label{tab:inv-p4-conceptos}
\end{table}

\section{Puntos de conexión con las preguntas}

\begin{itemize}
  \item \textbf{P4, inciso a):} la dependencia directa \texttt{OrderProcessor}
        $\rightarrow$ \texttt{AddressValidatorService} hace que una política de alto nivel
        nombre un detalle volátil, lo que contraviene el DIP \parencite[87]{martin2017} y la
        Regla de Dependencia \parencite[203]{martin2017}, y rompe la relación puerto-adaptador
        \parencite{cockburn2005}.
  \item \textbf{P4, inciso b):} el rediseño consiste en declarar \texttt{IAddressValidator} en
        el núcleo y dejar el adaptador HTTP fuera, en la capa de adaptadores de interfaz
        \parencite[205]{martin2017}, por ser el elemento volátil \parencite[88]{martin2017}.
  \item \textbf{P4, inciso c):} la prueba del núcleo sin HTTP se sostiene en la testabilidad de
        estas arquitecturas \parencite[202]{martin2017}, y el cambio de proveedor queda
        confinado al adaptador porque la dependencia apunta a la interfaz, el elemento más
        estable \parencite[88]{martin2017}.
  \item \textbf{P1, P2 y P6:} la Regla de Dependencia y la simetría puerto-adaptador son el
        mismo criterio que P1 y P2 contrastan con la N-Capas, y el argumento estructural que P6
        opone a las optimizaciones de capas abiertas.
\end{itemize}


 (ya presente en ese
% archivo: sin este fragmento, SemArq1_TeamBales.tex no compila).
% Fuente de trabajo: notas.md (misma carpeta).
% Citas: páginas impresas de Martin (2017); Cockburn (2005) y Palermo (2008) no
% tienen paginación y se citan por sección o parte. Todas las claves están en
% core.bib (martin2017, cockburn2005, palermo2008): fuentes.bib no declara nada.
% Comillas: se usan comillas tipográficas literales, no \enquote, para no
% depender de csquotes (no está en la lista de paquetes de preamble.tex, #13).
% =============================================================================

% Grupo local: \emergencystretch evita líneas desbordadas sin afectar a otros capítulos.
\begingroup\emergencystretch=3em \hbadness=2000
\section*{Notas de investigación: DIP, Inversión de Dependencias y diseño de interfaces en Clean Architecture}

\section{Síntesis}

El Principio de Inversión de Dependencias (DIP) sostiene que los sistemas más flexibles son
aquellos en los que las dependencias del código fuente se refieren solo a abstracciones, no a
concreciones \parencite[87]{martin2017}. La consecuencia práctica es que una política de alto
nivel —una regla de negocio— no debe nombrar un detalle de bajo nivel: ni una base de datos, ni
un servicio externo, ni una librería de transporte. Ambos, política y detalle, dependen de una
abstracción declarada en el núcleo. Martin añade el criterio que convierte el principio en
decisión de diseño: toda modificación de una interfaz abstracta obliga a modificar sus
implementaciones, pero no al contrario, de modo que las interfaces son menos volátiles que las
implementaciones \parencite[88]{martin2017}; de ahí las reglas operativas de no referirse a
clases concretas volátiles, no sobreescribir funciones concretas y recurrir a factorías
abstractas para la creación de objetos \parencite[89]{martin2017}.

La Regla de Dependencia formaliza el principio a escala arquitectónica: las dependencias del
código fuente deben apuntar solo hacia adentro, hacia políticas de mayor nivel, y nada de un
círculo interno puede conocer nada de un círculo externo \parencite[203]{martin2017}. Martin
presenta esa regla como denominador común de un conjunto de arquitecturas —Hexagonal de
Cockburn, DCI, BCE— que persiguen la misma separación de intereses y producen sistemas
independientes de frameworks, de la interfaz de usuario, de la base de datos y de cualquier
agencia externa, y con reglas de negocio probables sin ninguno de esos elementos
\parencite[202]{martin2017}. En la Hexagonal, el núcleo declara los \emph{puertos} y la
tecnología aporta los \emph{adaptadores} que los implementan \parencite{cockburn2005}; en la
Onion, el Modelo de Dominio ocupa el centro y las interfaces de repositorio viven en el núcleo
mientras el comportamiento que toca la base de datos queda fuera \parencite{palermo2008}.

El mecanismo concreto es siempre el mismo: el detalle implementa la abstracción que el núcleo
declara, y la traducción entre el contrato del núcleo y la tecnología externa ocurre en la capa
de adaptadores de interfaz \parencite[205]{martin2017}. Conviene no confundirlo con la inversión
del flujo de control: el núcleo sigue llamando a la validación o a la consulta; lo que se
invierte es la dirección de la dependencia del código fuente, que en el límite cruzado pasa a
apuntar desde el detalle hacia la abstracción \parencite[87, 91]{martin2017}. Ese es el punto de
apoyo de la testabilidad: con el puerto en el núcleo, la prueba sustituye el adaptador real por
un doble controlado y ejercita la regla de negocio sin red, sin servidor y sin base de datos
\parencite[202]{martin2017}.

\section{Conceptos clave}

\begin{table}[htbp]
  \centering
  \footnotesize
  \renewcommand{\arraystretch}{1.2}
  \begin{tabular}{@{}p{0.22\linewidth}p{0.46\linewidth}p{0.24\linewidth}@{}}
    \toprule
    \textbf{Concepto} & \textbf{Definición} & \textbf{Fuente} \\
    \midrule
    \raggedright Principio de Inversión de Dependencias (DIP)
      & \raggedright Los sistemas más flexibles son los que dependen solo de abstracciones, no
        de concreciones; la política de alto nivel no nombra el detalle
      & \raggedright \textcite[87]{martin2017} \tabularnewline
    \raggedright Volatilidad de interfaces frente a implementaciones
      & \raggedright Cambiar una interfaz obliga a cambiar sus implementaciones, pero no al
        contrario: la interfaz es el elemento más estable del que depender
      & \raggedright \textcite[88]{martin2017} \tabularnewline
    \raggedright Reglas operativas del DIP
      & \raggedright No referirse a clases concretas volátiles, no sobreescribir funciones
        concretas y usar factorías abstractas para la creación
      & \raggedright \textcite[89]{martin2017} \tabularnewline
    \raggedright Regla de Dependencia
      & \raggedright Las dependencias del código fuente apuntan solo hacia adentro; nada de un
        círculo interno conoce nada de un círculo externo
      & \raggedright \textcite[203]{martin2017} \tabularnewline
    \raggedright Características comunes de Clean, Hexagonal, DCI y BCE
      & \raggedright Independencia de frameworks, de la interfaz de usuario, de la base de datos
        y de agencias externas; reglas de negocio probables en aislamiento
      & \raggedright \textcite[202]{martin2017} \tabularnewline
    \raggedright Puertos y adaptadores
      & \raggedright El núcleo declara el puerto y la tecnología aporta el adaptador que lo
        implementa, de modo que el núcleo ignora el dispositivo concreto
      & \raggedright \textcite{cockburn2005} \tabularnewline
    \raggedright Capa de adaptadores de interfaz
      & \raggedright Convierte los datos entre el formato del caso de uso y el de la tecnología
        externa; es el lugar natural del detalle
      & \raggedright \textcite[205]{martin2017} \tabularnewline
    \raggedright Dominio al centro (Onion)
      & \raggedright El Modelo de Dominio no se acopla más que a sí mismo; las interfaces de
        repositorio van en el núcleo y el acceso a datos fuera
      & \raggedright \textcite{palermo2008} \tabularnewline
    \raggedright Flujo de control frente a dependencia de código fuente
      & \raggedright El núcleo sigue invocando la operación; lo que se invierte en el límite es
        la dirección de la dependencia del código fuente
      & \raggedright \textcite[87, 91]{martin2017} \tabularnewline
    \raggedright Testabilidad mediante dobles de prueba
      & \raggedright Con el puerto en el núcleo, la prueba sustituye el adaptador real y ejercita
        la regla de negocio sin interfaz, base de datos ni servidor
      & \raggedright \textcite[202]{martin2017} \tabularnewline
    \bottomrule
  \end{tabular}
  \caption{Conceptos clave de la investigación INV-P4.}
  \label{tab:inv-p4-conceptos}
\end{table}

\section{Puntos de conexión con las preguntas}

\begin{itemize}
  \item \textbf{P4, inciso a):} la dependencia directa \texttt{OrderProcessor}
        $\rightarrow$ \texttt{AddressValidatorService} hace que una política de alto nivel
        nombre un detalle volátil, lo que contraviene el DIP \parencite[87]{martin2017} y la
        Regla de Dependencia \parencite[203]{martin2017}, y rompe la relación puerto-adaptador
        \parencite{cockburn2005}.
  \item \textbf{P4, inciso b):} el rediseño consiste en declarar \texttt{IAddressValidator} en
        el núcleo y dejar el adaptador HTTP fuera, en la capa de adaptadores de interfaz
        \parencite[205]{martin2017}, por ser el elemento volátil \parencite[88]{martin2017}.
  \item \textbf{P4, inciso c):} la prueba del núcleo sin HTTP se sostiene en la testabilidad de
        estas arquitecturas \parencite[202]{martin2017}, y el cambio de proveedor queda
        confinado al adaptador porque la dependencia apunta a la interfaz, el elemento más
        estable \parencite[88]{martin2017}.
  \item \textbf{P1, P2 y P6:} la Regla de Dependencia y la simetría puerto-adaptador son el
        mismo criterio que P1 y P2 contrastan con la N-Capas, y el argumento estructural que P6
        opone a las optimizaciones de capas abiertas.
\end{itemize}


 (ya presente en ese
% archivo: sin este fragmento, SemArq1_TeamBales.tex no compila).
% Fuente de trabajo: notas.md (misma carpeta).
% Citas: páginas impresas de Martin (2017); Cockburn (2005) y Palermo (2008) no
% tienen paginación y se citan por sección o parte. Todas las claves están en
% core.bib (martin2017, cockburn2005, palermo2008): fuentes.bib no declara nada.
% Comillas: se usan comillas tipográficas literales, no \enquote, para no
% depender de csquotes (no está en la lista de paquetes de preamble.tex, #13).
% =============================================================================

% Grupo local: \emergencystretch evita líneas desbordadas sin afectar a otros capítulos.
\begingroup\emergencystretch=3em \hbadness=2000
\section*{Notas de investigación: DIP, Inversión de Dependencias y diseño de interfaces en Clean Architecture}

\section{Síntesis}

El Principio de Inversión de Dependencias (DIP) sostiene que los sistemas más flexibles son
aquellos en los que las dependencias del código fuente se refieren solo a abstracciones, no a
concreciones \parencite[87]{martin2017}. La consecuencia práctica es que una política de alto
nivel —una regla de negocio— no debe nombrar un detalle de bajo nivel: ni una base de datos, ni
un servicio externo, ni una librería de transporte. Ambos, política y detalle, dependen de una
abstracción declarada en el núcleo. Martin añade el criterio que convierte el principio en
decisión de diseño: toda modificación de una interfaz abstracta obliga a modificar sus
implementaciones, pero no al contrario, de modo que las interfaces son menos volátiles que las
implementaciones \parencite[88]{martin2017}; de ahí las reglas operativas de no referirse a
clases concretas volátiles, no sobreescribir funciones concretas y recurrir a factorías
abstractas para la creación de objetos \parencite[89]{martin2017}.

La Regla de Dependencia formaliza el principio a escala arquitectónica: las dependencias del
código fuente deben apuntar solo hacia adentro, hacia políticas de mayor nivel, y nada de un
círculo interno puede conocer nada de un círculo externo \parencite[203]{martin2017}. Martin
presenta esa regla como denominador común de un conjunto de arquitecturas —Hexagonal de
Cockburn, DCI, BCE— que persiguen la misma separación de intereses y producen sistemas
independientes de frameworks, de la interfaz de usuario, de la base de datos y de cualquier
agencia externa, y con reglas de negocio probables sin ninguno de esos elementos
\parencite[202]{martin2017}. En la Hexagonal, el núcleo declara los \emph{puertos} y la
tecnología aporta los \emph{adaptadores} que los implementan \parencite{cockburn2005}; en la
Onion, el Modelo de Dominio ocupa el centro y las interfaces de repositorio viven en el núcleo
mientras el comportamiento que toca la base de datos queda fuera \parencite{palermo2008}.

El mecanismo concreto es siempre el mismo: el detalle implementa la abstracción que el núcleo
declara, y la traducción entre el contrato del núcleo y la tecnología externa ocurre en la capa
de adaptadores de interfaz \parencite[205]{martin2017}. Conviene no confundirlo con la inversión
del flujo de control: el núcleo sigue llamando a la validación o a la consulta; lo que se
invierte es la dirección de la dependencia del código fuente, que en el límite cruzado pasa a
apuntar desde el detalle hacia la abstracción \parencite[87, 91]{martin2017}. Ese es el punto de
apoyo de la testabilidad: con el puerto en el núcleo, la prueba sustituye el adaptador real por
un doble controlado y ejercita la regla de negocio sin red, sin servidor y sin base de datos
\parencite[202]{martin2017}.

\section{Conceptos clave}

\begin{table}[htbp]
  \centering
  \footnotesize
  \renewcommand{\arraystretch}{1.2}
  \begin{tabular}{@{}p{0.22\linewidth}p{0.46\linewidth}p{0.24\linewidth}@{}}
    \toprule
    \textbf{Concepto} & \textbf{Definición} & \textbf{Fuente} \\
    \midrule
    \raggedright Principio de Inversión de Dependencias (DIP)
      & \raggedright Los sistemas más flexibles son los que dependen solo de abstracciones, no
        de concreciones; la política de alto nivel no nombra el detalle
      & \raggedright \textcite[87]{martin2017} \tabularnewline
    \raggedright Volatilidad de interfaces frente a implementaciones
      & \raggedright Cambiar una interfaz obliga a cambiar sus implementaciones, pero no al
        contrario: la interfaz es el elemento más estable del que depender
      & \raggedright \textcite[88]{martin2017} \tabularnewline
    \raggedright Reglas operativas del DIP
      & \raggedright No referirse a clases concretas volátiles, no sobreescribir funciones
        concretas y usar factorías abstractas para la creación
      & \raggedright \textcite[89]{martin2017} \tabularnewline
    \raggedright Regla de Dependencia
      & \raggedright Las dependencias del código fuente apuntan solo hacia adentro; nada de un
        círculo interno conoce nada de un círculo externo
      & \raggedright \textcite[203]{martin2017} \tabularnewline
    \raggedright Características comunes de Clean, Hexagonal, DCI y BCE
      & \raggedright Independencia de frameworks, de la interfaz de usuario, de la base de datos
        y de agencias externas; reglas de negocio probables en aislamiento
      & \raggedright \textcite[202]{martin2017} \tabularnewline
    \raggedright Puertos y adaptadores
      & \raggedright El núcleo declara el puerto y la tecnología aporta el adaptador que lo
        implementa, de modo que el núcleo ignora el dispositivo concreto
      & \raggedright \textcite{cockburn2005} \tabularnewline
    \raggedright Capa de adaptadores de interfaz
      & \raggedright Convierte los datos entre el formato del caso de uso y el de la tecnología
        externa; es el lugar natural del detalle
      & \raggedright \textcite[205]{martin2017} \tabularnewline
    \raggedright Dominio al centro (Onion)
      & \raggedright El Modelo de Dominio no se acopla más que a sí mismo; las interfaces de
        repositorio van en el núcleo y el acceso a datos fuera
      & \raggedright \textcite{palermo2008} \tabularnewline
    \raggedright Flujo de control frente a dependencia de código fuente
      & \raggedright El núcleo sigue invocando la operación; lo que se invierte en el límite es
        la dirección de la dependencia del código fuente
      & \raggedright \textcite[87, 91]{martin2017} \tabularnewline
    \raggedright Testabilidad mediante dobles de prueba
      & \raggedright Con el puerto en el núcleo, la prueba sustituye el adaptador real y ejercita
        la regla de negocio sin interfaz, base de datos ni servidor
      & \raggedright \textcite[202]{martin2017} \tabularnewline
    \bottomrule
  \end{tabular}
  \caption{Conceptos clave de la investigación INV-P4.}
  \label{tab:inv-p4-conceptos}
\end{table}

\section{Puntos de conexión con las preguntas}

\begin{itemize}
  \item \textbf{P4, inciso a):} la dependencia directa \texttt{OrderProcessor}
        $\rightarrow$ \texttt{AddressValidatorService} hace que una política de alto nivel
        nombre un detalle volátil, lo que contraviene el DIP \parencite[87]{martin2017} y la
        Regla de Dependencia \parencite[203]{martin2017}, y rompe la relación puerto-adaptador
        \parencite{cockburn2005}.
  \item \textbf{P4, inciso b):} el rediseño consiste en declarar \texttt{IAddressValidator} en
        el núcleo y dejar el adaptador HTTP fuera, en la capa de adaptadores de interfaz
        \parencite[205]{martin2017}, por ser el elemento volátil \parencite[88]{martin2017}.
  \item \textbf{P4, inciso c):} la prueba del núcleo sin HTTP se sostiene en la testabilidad de
        estas arquitecturas \parencite[202]{martin2017}, y el cambio de proveedor queda
        confinado al adaptador porque la dependencia apunta a la interfaz, el elemento más
        estable \parencite[88]{martin2017}.
  \item \textbf{P1, P2 y P6:} la Regla de Dependencia y la simetría puerto-adaptador son el
        mismo criterio que P1 y P2 contrastan con la N-Capas, y el argumento estructural que P6
        opone a las optimizaciones de capas abiertas.
\end{itemize}


